\documentclass[../../main]{subfiles}

\begin{document}

\section{Relações Binárias}
\label{sec:matematica:relacoes:relacoes-binarias}

\begin{defn}[Relação Binária]
  \label{def:matematica:relacoes:relacoes-binarias:relacao-binaria}

  Sejam \(A\) e \(B\) conjuntos.

  Uma relação binária de \(A\) em \(B\) é qualquer subconjunto de

  \[
    A \times B.
  \]

  Isto é,

  \[
    R \subseteq A \times B.
  \]

\end{defn}

\begin{defn}[Elementos Relacionados]
  \label{def:matematica:relacoes:relacoes-binarias:elementos-relacionados}

  Seja \(R\subseteq A\times B\).

  Dizemos que \(a\in A\) está relacionado com \(b\in B\) por \(R\)
  quando

  \[
    (a,b)\in R.
  \]

  Neste caso escrevemos

  \[
    aRb.
  \]

\end{defn}

\begin{defn}[Domínio de uma Relação]
  \label{def:matematica:relacoes:relacoes-binarias:dominio}

  Seja \(R\subseteq A\times B\).

  O domínio de \(R\) é o conjunto

  \[
    \dom(R)
    =
    \{
      a\in A
      \mid
      (\exists b\in B)
      ((a,b)\in R)
    \}.
  \]

\end{defn}

\begin{defn}[Imagem de uma Relação]
  \label{def:matematica:relacoes:relacoes-binarias:imagem}

  Seja \(R\subseteq A\times B\).

  A imagem de \(R\) é o conjunto

  \[
    \operatorname{Im}(R)
    =
    \{
      b\in B
      \mid
      (\exists a\in A)
      ((a,b)\in R)
    \}.
  \]

\end{defn}

\begin{defn}[Relação Identidade]
  \label{def:matematica:relacoes:relacoes-binarias:identidade}

  Seja \(A\) um conjunto.

  A relação identidade em \(A\) é definida por

  \[
    \Delta_A
    =
    \{
      (a,a)
      \mid
      a\in A
    \}.
  \]

\end{defn}

\begin{defn}[Relação Inversa]
  \label{def:matematica:relacoes:relacoes-binarias:inversa}

  Seja \(R\subseteq A\times B\).

  A relação inversa de \(R\) é definida por

  \[
    R^{-1}
    =
    \{
      (b,a)
      \mid
      (a,b)\in R
    \}.
  \]

\end{defn}

\begin{prop}
  \label{prop:matematica:relacoes:relacoes-binarias:inversa-inversa}

  Para toda relação \(R\),

  \[
    (R^{-1})^{-1}
    =
    R.
  \]

\end{prop}

\begin{proof}

  Segue diretamente da definição de relação inversa.

\end{proof}

\end{document}
